\section{The Capacity-Grid List Theorem}

The following argument applies to every injectively indexed Reed--Solomon
evaluation domain; smoothness enters only when we instantiate the explicit
family from \cref{sec:explicit-family}.

\begin{theorem}[Universal capacity-grid equality]
\label{thm:capacity-grid-exact}
Let \(\mathcal C=\RS[\F,\mathcal L,k]\) have length \(n\), where
\(1\le k\le n\), and let \(m\ge1\).  Then
\[
  L_m\left(\mathcal C,\frac{n-k}{n}\right)=\binom nk .
\]
In particular, the worst-case endpoint list size is independent of \(m\).
\end{theorem}

\begin{proof}
We first prove the upper bound.  Fix \(w\in(\F^m)^n\).
Every codeword in the ball agrees with \(w\) on at least \(k\) complete
vector-symbol coordinates.  Fix a total order on the \(k\)-subsets of \(I\)
and assign to each codeword its first \(k\)-subset of agreement coordinates.
Two interleaved codewords assigned the same set \(S\) agree with one another
on \(S\).  In each component, two degree-\(<k\) polynomials agree at the \(k\)
distinct evaluation points indexed by \(S\), and hence are equal.  The two
interleaved codewords are therefore equal, so the assignment is injective
into the \(\binom nk\) possible sets \(S\).

For the reverse inequality, choose one of the \(m\) components and define the
received word
\[
  w_i^\star=(a_i^k,0,\ldots,0).
\]
For \(k<n\), the degree-\(k\) component is a classical RS deep-hole
witness \cite[Section~2]{ketiWan2016deepHoles}. The scalar list is the
empty-prefix case of Chojecki's exact prefix-list bijection
\cite[Section~4]{chojecki2026thresholds}. The interpolation argument below
counts its full interleaved endpoint list explicitly.
For each \(k\)-subset \(S\subseteq I\), put
\[
  V_S(X)=\prod_{i\in S}(X-a_i),
  \qquad
  f_S(X)=X^k-V_S(X).
\]
The polynomial \(V_S\) is monic of degree \(k\), so its leading term cancels
that of \(X^k\) and \(\deg f_S<k\).  Use \(f_S\) in the distinguished
component and zero polynomials in all other components to obtain an
interleaved codeword \(c_S\).

At evaluation coordinate \(i\),
\[
  c_{S,i}=w_i^\star
  \quad\Longleftrightarrow\quad
  V_S(a_i)=0
  \quad\Longleftrightarrow\quad
  i\in S,
\]
where the final equivalence uses injectivity of the evaluation points.
Consequently \(c_S\) has Hamming count distance exactly \(n-k\) from
\(w^\star\), and its agreement set recovers \(S\).  The \(\binom nk\) choices
of \(S\) give distinct codewords in the endpoint ball, matching the upper bound.
\end{proof}

\begin{theorem}[Failure on the next Hamming shell]
\label{thm:next-shell-lower}
Under the same hypotheses,
\[
  \left|
    \List\left(
      \mathcal C^{\equiv m},
      0,
      \frac{n-k+1}{n}
    \right)
  \right|
  \ge |\F|^m .
\]
\end{theorem}

\begin{proof}
Fix \(k-1\) evaluation points \(S\subseteq\mathcal L\) and set
\[
  g_S(X)=\prod_{a\in S}(X-a).
\]
For every \(\lambda=(\lambda_1,\ldots,\lambda_m)\in\F^m\), the tuple
\[
  (\lambda_1g_S,\ldots,\lambda_mg_S)
\]
defines an interleaved codeword.  These \(|\F|^m\) codewords are distinct.
The zero tuple has distance zero from the received word \(0\). Every nonzero
tuple vanishes on \(S\) and is nonzero at every evaluation point outside
\(S\), so its Hamming count distance is exactly \(n-k+1\).
\end{proof}

\begin{theorem}[Fractional-target phase split]
\label{thm:fractional-phase-split}
Fix \(0<\epsstar<1\), write \(q=|\F|\), and call a radius \(\delta\)
safe when
\[
  L_m(\mathcal C,\delta)\le\epsstar q.
\]
For every \(m\ge1\), exactly one of the following regimes applies:
\begin{enumerate}[label=(\roman*)]
  \item If \(\epsstar q<1\), no radius in \([0,1]\) is safe.
  \item If
    \(1\le\epsstar q<\binom nk\), radius zero is safe but the capacity
    endpoint and every larger radius fail.
  \item If \(\binom nk\le\epsstar q\), the largest safe radius on
    \(\mathcal H_n\) is
    \[
      \deltacap=\frac{n-k}{n}=1-\rho;
    \]
    and over real radii the safe set is exactly
    \[
      \left[0,\frac{n-k+1}{n}\right),
    \]
    so it has no largest element and has supremum
    \((n-k+1)/n\).
 \end{enumerate}
Equivalently, the fractional target holds at the capacity endpoint if and
only if
\[
  \binom nk\le\epsstar|\F|.
\]
\end{theorem}

\begin{proof}
Every codeword is contained in the radius-zero ball centered at itself, so no
radius can be safe when \(\epsstar q<1\).  A radius-zero ball contains at most
one codeword, proving its safety when \(\epsstar q\ge1\).

\Cref{thm:capacity-grid-exact} proves both the failure in the intermediate
regime and the necessary-and-sufficient endpoint criterion.  In the final
regime, it and monotonicity of list balls prove safety through count radius
\(n-k\).  A real radius strictly below
\((n-k+1)/n\) defines the same ball as some count radius at most \(n-k\).
At the next grid point, \cref{thm:next-shell-lower} gives a list of size at
least \(|\F|^m\ge|\F|>\epsstar|\F|\); monotonicity then rules out every
larger radius.
\end{proof}

\begin{corollary}[The explicit smooth family]
\label{cor:explicit-family-list}
For every \(u\ge0\), every \(m\ge1\), and
\(\mathcal C_u=\RS[\F_u,H_u,k_u]\) from
\cref{sec:explicit-family}, the conclusions of
\cref{thm:fractional-phase-split}\textup{(iii)} hold with
\(\epsstar=2^{-128}\).  In
particular, the largest safe Hamming-grid radius is \(3/4\).
\end{corollary}

\begin{proof}
Apply the inequality proved in \cref{sec:field-inequality}.
\end{proof}

\subsection{The one-codeword budget}

The high-agreement ledger in \cite{chojecki2026results} records the same
unique-list radius for all interleaving widths. We give the minimum-distance
argument and an explicit two-codeword witness to pin the first unsafe count.

\begin{theorem}[Exact list threshold for a one-codeword budget]
\label{thm:one-codeword-budget}
Let \(1\le k\le n\), \(m\ge1\), and \(1\le\epsstar q<2\). Set
\(d=n-k+1\). For every natural count radius \(r\le n\),
\[
  L_m(\mathcal C,r/n)\le\epsstar q
  \quad\Longleftrightarrow\quad
  2r<d.
\]
The last safe count is \(\lfloor(d-1)/2\rfloor\), the first unsafe count
is \(\lceil d/2\rceil\), and the real safe set is
\([0,\lceil d/2\rceil/n)\).
\end{theorem}

\begin{proof}
An integer list size lies within the budget exactly when it is at most one.
Any two distinct interleaved codewords differ by a nonzero degree-\(<k\)
polynomial in at least one component. The root bound leaves at least
\(d\) coordinates where that component differs, so the vector-symbol
Hamming distance is at least \(d\). If both codewords were in the same
count-\(r\) ball, the triangle inequality would give distance at most
\(2r\). Thus every ball has size at most one when \(2r<d\).

The zero codeword and the evaluation of \(g_S\) from
\cref{thm:next-shell-lower} in one component, with zeros in the others,
differ at exactly \(d\) sites. Split these sites into two sets of sizes
\(\lfloor d/2\rfloor\) and \(\lceil d/2\rceil\). On the first set let a
received word take the nonzero codeword's symbols, and on all other sites
take the zero codeword's symbols. Its distances to the two codewords
are the two set sizes. Its ball therefore contains both codewords at
count \(\lceil d/2\rceil\) and at every larger count. Every relative
distance is an integer multiple of \(1/n\). A real radius strictly below
\(\lceil d/2\rceil/n\) thus admits only counts at most
\(\lceil d/2\rceil-1\), proving the real-radius claim.
\end{proof}

\begin{corollary}[The boundary prime field]
\label{cor:boundary-field-list}
For the smooth code over \(q=2^{128}+192513\) with \(n=4096\) and
\(k=1024\) from \cref{sec:mca-boundary-field}, the \(2^{-128}\) list target
has last safe count \(1536\), first unsafe count \(1537\), and real safe
set \([0,1537/4096)\) for every positive interleaving width. For the same
code, the corresponding pair-MCA counts are \(0\) and \(1\) by
\cref{thm:mca-boundary-threshold}.
\end{corollary}

\begin{proof}
Here \(1\le2^{-128}q<2\) and \(d=3073\). Apply
\cref{thm:one-codeword-budget}.
\end{proof}

\subsection{Packing agreement subsets below capacity}

\begin{theorem}[Finite-radius agreement packing]
\label{thm:list-agreement-packing}
Let \(\mathcal C=\RS[\F,\mathcal L,k]\) have length \(n\), where
\(1\le k<n\), and let \(m\ge1\). For \(0\le r\le n-k\),
\[
  L_m(\mathcal C,r/n)\le
  \left\lfloor\frac{\binom nk}{\binom{n-r}k}\right\rfloor.
\]
In particular, if the integer list budget \(B\) satisfies
\[
  \left\lfloor\frac{\binom nk}{k+1}\right\rfloor\le B<\binom nk,
\]
the last safe count is \(n-k-1\), the first unsafe count is \(n-k\),
and the real safe set is \([0,(n-k)/n)\).
\end{theorem}

\begin{proof}
Fix a received word and give each candidate in its count-\(r\) list every
\(k\)-subset of its agreement coordinates. Each candidate receives at least
\(\binom{n-r}{k}\) subsets. If two candidates received the same subset,
interpolation in every component would make them equal. The candidates'
subset collections are disjoint, so their total size is at most \(\binom nk\).
At count \(n-k-1\), the denominator is \(\binom{k+1}k=k+1\).
At the next count, \cref{thm:capacity-grid-exact} attains the upper endpoint
\(\binom nk\). Monotonicity and the half-open Hamming shells give the
stated thresholds.
\end{proof}

\subsection{Exact paired thresholds on four length-64 domains}
\label{sec:jo-paired-thresholds}

Use the four length-\(64\) codes in Tables~2--3 of Jo
\cite[Section~5.2]{jo2026beyondJohnson}. Each domain is an order-\(64\)
subgroup of the corresponding prime field. The displayed MCA counts are
Jo's results. Our list counts use the same fields, code dimensions, and
evaluation domains. Each pair in the table gives the last safe and first
unsafe count at \(\epsstar=2^{-128}\).

\begin{center}
\begin{tabular}{@{}rrrr@{}}
\toprule
Rate & List budget \(\lfloor 2^{-128}q\rfloor\) & List counts & Pair-MCA counts\\
\midrule
\(1/2\)  & \(2^{58}-1\) & \((31,32)\) & \((30,31)\)\\
\(1/4\)  & \(2^{49}-1\) & \((48,49)\) & \((46,47)\)\\
\(1/8\)  & \(2^{33}-1\) & \((56,57)\) & \((54,55)\)\\
\(1/16\) & \(2^{22}-1\) & \((60,61)\) & \((58,59)\)\\
\bottomrule
\end{tabular}
\end{center}

For \(k=32\), the budget lies strictly below \(\binom{64}{32}\) and
above \(\lfloor\binom{64}{32}/33\rfloor=55\,534\,064\,877\,048\,198\).
\Cref{thm:list-agreement-packing} pins the list transition at capacity.
For \(k=16,8,4\), the respective budgets dominate \(\binom{64}{k}\),
so \cref{thm:fractional-phase-split} gives safety at capacity and failure
on the next shell. All four list statements hold at every positive
interleaving width. In each column the first unsafe count divided by
\(64\) is the excluded endpoint of the real safe interval.

\subsection{A subfield construction below capacity}
\label{sec:subfield-crossings}

Jo's local-test upper bound \cite[Theorem~4.2]{jo2026beyondJohnson} and
Chojecki's locator-prefix and pole constructions
\cite[Section~4]{chojecki2026thresholds} can pin paired thresholds below
capacity. The following conditions use only integer comparisons.

\begin{proposition}[Paired threshold certificate from a subfield]
\label{prop:subfield-crossings}
Let the \(n\) evaluation points of \(\mathcal C=\RS[\mathbb F_q,\mathcal L,k]\)
lie in a subfield \(\mathbb F_b\), with \(1\le k\le n-3\). Set
\(B=\lfloor 2^{-128}q\rfloor\). Assume
\begin{align*}
 \max\left\{\left\lfloor\frac{\binom nk}{\binom{k+2}k}\right\rfloor,
 \left\lfloor\frac{\binom n{k+1}}{\binom{k+2}k}\right\rfloor\right\}
     &\le B,\\
 B&<\min\left\{\left\lceil\frac{\binom n{k+1}}b\right\rceil,
 \left\lceil\frac{\binom n{k+2}}b\right\rceil\right\},\\
 n+k\binom{B+1}{2}&<q.
\end{align*}
Then the 128-bit list counts are \((n-k-2,n-k-1)\) at every positive
interleaving width, and the pair-MCA counts are \((n-k-3,n-k-2)\), with
last safe count first in each pair.
\end{proposition}

\begin{proof}
For the list safe count, apply \cref{thm:list-agreement-packing} at
agreement size \(k+2\). For the MCA safe count, Jo's integer bound at
agreement size \(k+3\) is
\(\lfloor\binom n{k+1}/\binom{k+2}k\rfloor\).

For each \(a\in\{k+1,k+2\}\), partition the \(a\)-subsets \(S\) of
\(\mathcal L\) by their sum \(s\in\mathbb F_b\). Some class contains at
least \(\lceil\binom na/b\rceil\ge B+1\) sets. In that class put
\[
 U(X)=X^a-sX^{a-1},\qquad
 P_S(X)=U(X)-\prod_{x\in S}(X-x).
\]
The two leading coefficients cancel, so \(\deg P_S\le a-2\). Distinct
sets give distinct \(P_S\), and \(P_S\) agrees with \(U\) on \(S\).
For \(a=k+1\), evaluation preserves distinctness by the root bound.
These are \(B+1\) distinct list entries at count
\(n-k-1\). Placing them in one component and filling the other components
with zero gives the same witness at every positive width.

For \(a=k+2\), choose \(B+1\) of the \(P_S\). Each difference has degree
at most \(k\), so at most \(k\binom{B+1}{2}\) poles cause any equality of
their evaluations. The field-size inequality gives a pole
\(\alpha\in\mathbb F_q\setminus\mathcal L\) where all their values are
distinct. Use the affine line of received words with symbols
\[
 y_\gamma(x)=\frac{U(x)-\gamma}{x-\alpha}
 \qquad(x\in\mathcal L).
\]
At \(\gamma=P_S(\alpha)\), the polynomial
\[
 \frac{P_S(X)-P_S(\alpha)}{X-\alpha}
\]
has degree at most \(k-1\) and agrees with \(y_\gamma\) on \(S\).
The direction \(-1/(x-\alpha)\) cannot agree on \(S\) with a polynomial
\(G\) of degree less than \(k\). Otherwise \((X-\alpha)G(X)+1\) would
have the \(k+2\) points of \(S\) as roots despite having degree at most
\(k\) and value one at \(\alpha\). Thus each distinct \(P_S(\alpha)\)
is a bad challenge on its own support at count \(n-k-2\). This gives
more than \(B\) bad challenges, which pins the MCA threshold.
\end{proof}

\begin{corollary}[An exact length-128 pair of thresholds]
\label{cor:subfield-128-thresholds}
Take the rate-\(1/2\) code on the order-\(128\) subgroup of
\(\mathbb F_{641}^{\times}\) inside \(\mathbb F_{641^{26}}\).
The list counts are \((62,63)\) at every positive width, and the pair-MCA
counts are \((61,62)\). The real safe sets are \([0,63/128)\) and
\([0,62/128)\), respectively.
\end{corollary}

\begin{proof}
The prime \(641\) has \(128\mid640\), giving the smooth subgroup, and
\(q=641^{26}<2^{256}\). The budget is
\(B=27955632488535336577577692993784094\).
The two safe-side bounds from \cref{prop:subfield-crossings} are
\(\lfloor\binom{128}{64}/2145\rfloor\) and
\(\lfloor\binom{128}{65}/2145\rfloor\), both at most \(B\).
The two prefix-class bounds \(\lceil\binom{128}{65}/641\rceil\) and
\(\lceil\binom{128}{66}/641\rceil\) both exceed \(B\), and
\(128+64\binom{B+1}{2}<641^{26}\). The proposition gives the count
thresholds and the next Hamming shell gives each real endpoint.
\end{proof}

\subsection{List sizes under extension of the domain field}
\label{sec:list-descent}

Let \(D\subset K=\mathbb F_b\) be the fixed evaluation domain and let
\(F=\mathbb F_{b^e}\). For a multiplicative coset of order \(n\), rescaling
the sites to the subgroup preserves evaluations after a change of variable.
In characteristic \(p\), the subgroup generates the field of degree
\(\operatorname{ord}_n(p)\) over \(\mathbb F_p\) when \(n>1\).
Expanding every symbol and every polynomial
coefficient in a \(K\)-basis of \(F\) identifies width \(m\) over \(F\)
with width \(me\) over \(K\). Polynomial evaluation at the sites in \(D\)
commutes with this expansion. The entire vector symbol at a site agrees
exactly when all its basis coordinates agree, so the identification
preserves each count ball and its list size. The extension-code arrangement
and its relation to list weight enumerators also appear in
\cite{jurriusPellikaan2015}.
For a center already defined over \(K\) and \(r\le n-k\), every candidate
agrees at \(k\) sites with values in \(K\), so interpolation puts its
coefficients in \(K\) as well.

For \(r\le R=n-k\), use a parity check over \(K\) with syndrome space
\(V=K^R\). Write \(U_E\) for the span of the columns on support \(E\).
Let \(W\le V\) be the span of the \(me\) scalar syndrome vectors in an
expanded word. This span has an intrinsic description. Its extension to
\(F\) is the \(F\)-span of the Galois conjugates of the original \(m\)
syndrome columns. Each conjugate is a linear combination of the basis
coefficients. Conversely, the trace-dual basis recovers each coefficient
vector as the trace of a scalar multiple of its original column, a sum of
weighted conjugates.

The list size of the word is exactly
\begin{equation}
 \ell_r(W)=\#\{E:|E|\le r,\ W\subseteq U_E,\
         W\nsubseteq U_{E\setminus\{i\}}\text{ for every }i\in E\}.
 \label{eq:list-exact-support-descent}
\end{equation}
The columns on \(E\) are independent because \(r\le R\), so they give
one error vector on that support for each expanded syndrome in \(U_E\).
The error array has exact support \(E\) precisely when deleting any
coordinate would exclude some expanded syndrome. Every \(W\) of dimension
at most \(me\) is realized by an expanded syndrome tuple, since the parity
check is onto. The exact worst list is the maximum of
\eqref{eq:list-exact-support-descent} over these \(K\)-subspaces.

If two distinct list entries have supports \(E,E'\), their difference is
a nonzero codeword in at least one component, so \(|E\cup E'|\ge R+1\).
The two column spans together have rank \(R\). It follows that
\(\dim W\le\dim(U_E\cap U_{E'})=|E|+|E'|-R\le2r-R\).
Consequently the worst list size stabilizes once
\(me\ge\max(1,2r-R)\), with value one at every width if no list has
two entries. The stabilization concerns list sizes at a fixed count.
The fractional budget still depends on \(|F|\).

One can compute the maximum using the intersection lattice of the \(U_E\),
including \(V\). For a span \(W\), intersect all the \(U_E\) containing
it to obtain its smallest lattice element \(X\). The membership tests in
\eqref{eq:list-exact-support-descent} depend only on \(X\). The number
\(N(X)\) of expanded syndrome tuples with this smallest element satisfies
\[
 (b^{me})^{\dim X}=\sum_{Y\subseteq X}N(Y),
\]
where the sum runs over lattice elements. Subtracting counts for smaller
elements gives \(N(X)\). The maximum is \(\max_{N(X)>0}\ell_r(X)\).
The computation still requires the intersection lattice for the supplied
domain.

\subsection{Exact thresholds from support ranks}
\label{sec:list-support-ranks}

Write \(R=n-k\), \(q=|\F|\), and let \(B\ge1\) be an integer list
budget. For \(0\le r\le R\), let \(\mathcal S_r\) be the supports of size
at most \(r\), and let \(U_E\) be the parity-column span on \(E\). All
support spaces in this subsection are over \(\F\).

\begin{theorem}[Support-rank threshold test]
\label{thm:list-support-ranks}
Suppose \((B+1)r<q\). At every positive interleaving width the list budget
is exceeded at count \(r\) if and only if there are \(B+1\) distinct
supports \(E_a\in\mathcal S_r\) whose intersection \(X=\bigcap_a U_{E_a}\)
satisfies
\[
 X\nsubseteq U_{E_a\setminus\{i\}} \qquad (i\in E_a).
\]
Equivalently, when \(2r>R\), the budget is exceeded if and only if
\(\ell_r(X)>B\) for some nonzero intersection \(X\) of at most
\(2r-R+1\) support spaces. Here \(\ell_r\) is defined by
\eqref{eq:list-exact-support-descent}. For \(2r\le R\) the budget holds.
\end{theorem}

\begin{proof}
A violating list of \(B+1\) codewords gives \(B+1\) errors with the same
syndrome tuple. Since the parity check is injective on each support of
size at most \(R\), the exact supports are distinct. Each syndrome
coordinate belongs to \(X\), and each site in an exact support has a nonzero
error in at least one component. Therefore some syndrome coordinate
lies outside the span obtained by deleting that site.

Conversely, each deletion cuts a proper linear subspace of \(X\).
If there is at least one deletion then \(\dim X\ge1\), and each proper
subspace has at most \(q^{\dim X-1}\) elements. There are at most
\((B+1)r<q\) deletions, so a scalar syndrome in \(X\) avoids all of them.
The parity check is onto. Choose a received word with this syndrome and
subtract the corresponding supported errors. This gives \(B+1\) distinct
codewords in the ball, also at every positive width by repeating the
scalar coordinate.

For a violating list, any two distinct exact supports \(E,E'\) have
\(\lvert E\cup E'\rvert\ge R+1\) by the minimum distance. The MDS column
property then gives \(\dim(U_E\cap U_{E'})\le 2r-R\). At least one
syndrome coordinate is nonzero, since the zero syndrome has only the
zero error on each such support. Start the intersection with this pair
and add another support only when it lowers the dimension. The dimension
stays positive, so at most \(2r-R+1\) supports define \(X\), which still
has \(\ell_r(X)>B\). In the reverse direction, choose \(B+1\) supports
counted by \(\ell_r(X)\). Their intersection contains \(X\), so it passes
the deletion test. When \(2r\le R\),
minimum distance alone gives at most one codeword.
\end{proof}

The test can use only ranks over the field generated by the evaluation
points. For each \(E\) form a matrix \(A_E\) whose rows are the coefficient
vectors of \(X^j\prod_{i\in E}(X-a_i)\), for \(0\le j<R-|E|\). These
rows span the annihilator of \(U_E\) for the weighted Vandermonde parity
check, since nonzero column weights leave \(U_E\) unchanged. For a support
intersection with stacked matrix \(A\), a support \(E\) contributes to
\(\ell_r\) exactly when stacking \(A_E\) leaves the rank unchanged and
stacking each \(A_{E\setminus\{i\}}\) increases it. Ranks are invariant
under extension of the domain field. If \(S=|\mathcal S_r|\) and
\(2r>R\), this gives an enumeration using at most
\(O(\operatorname{poly}(n) S^{2r-R+2})\) field operations. Its cost can
be exponential in the block length.

For \(B=\lfloor q/2^{128}\rfloor\ge1\) and \(n<2^{127}\),
\((B+1)r\le 2Br<2^{128}B\le q\) at every count \(r\le R\).
Thus the test determines the entire threshold from support ranks for
any such prescribed code. Take the first count with a violation, using
\(R+1\) if there is none below or at capacity, by the next-shell
construction. If that count is \(s\), the greatest safe count is
\(s-1\) and the real safe interval is \([0,s/n)\). A field with
\(B=0\) has no safe count. On a fixed domain over \(\mathbb F_b\), let
\(M_r\) be the maximum of one and the counts \(\ell_r(X)\) in the rank
test, with \(M_r=1\) for \(2r\le R\). The first extension degree at
which count \(r\le R\) is safe is the least \(e\ge1\) with
\(b^e\ge2^{128}M_r\). Lean proves the deletion criterion and its
128-bit specialization. The shortened rank enumeration is proved above.

\begin{corollary}[The two-codeword budget on half-rate subgroups]
\label{cor:list-two-half-rate}
Suppose the evaluation domain is a multiplicative coset of length
\(n=4t\) and \(k=2t\). At budget \(B=2\), the last safe and first unsafe
counts are \(t\) and \(t+1\), at every positive interleaving width.
The real safe interval is \([0,(t+1)/n)\).
\end{corollary}

\begin{proof}
The minimum distance is \(2t+1\), so count \(t\) is safe. The map
\(x\mapsto x^t\) partitions the coset into four fibers of size \(t\),
with distinct values \(b_0,b_1,b_2,b_3\). Choose a site \(z\) in the
fiber of \(b_0\), and set \(g(X)=(X^t-b_0)/(X-z)\). Define
\[
 L_1(T)=0,\qquad L_2(T)=T-b_3,\qquad
 L_3(T)=\frac{b_1-b_3}{b_1-b_2}(T-b_2).
\]
For \(i=1,2,3\), the two affine polynomials with indices other than
\(i\) agree at \(b_i\). The polynomials \(f_i(X)=g(X)L_i(X^t)\)
are distinct and have degree at most \(2t-1\), with distinct evaluations
by the degree bound. Set the received value on fiber \(i\) to this
common evaluation for \(i=1,2,3\), and set it to zero on fiber \(0\).
Each \(f_i\) can disagree only on fiber \(i\) and at \(z\), because
\(g\) vanishes at every other site of fiber \(0\). These three words
give a violating list at count \(t+1\). Repeating the scalar coordinate
gives the same list for each positive width.
\end{proof}

For power-of-two lengths \(n\ge16\), the lower list bound also follows
from Al Kindi's subgroup construction \cite[Theorem~1.1]{kindi2026nearCapacity}
using a subgroup of order \(n/8\). The four-fiber proof includes lengths
four and eight.

\subsection{Scope of the list-size guarantee}

\Cref{thm:capacity-grid-exact} completely determines the endpoint list size
for every supplied Reed--Solomon code, smooth or otherwise.  Thus the displayed
binomial field-size inequality is necessary and sufficient at that endpoint,
not merely a convenient sufficient hypothesis.
\Cref{cor:explicit-family-list} supplies a concrete smooth family in the safe
regime. \Cref{thm:one-codeword-budget} resolves every one-codeword budget.
\Cref{thm:list-agreement-packing} resolves the budgets immediately below
the capacity list size. \Cref{thm:list-support-ranks} determines each
prescribed-code threshold under its field-size condition. A general
closed-form threshold that avoids support enumeration remains open for
\(1\le k<n\) in the intermediate budget range
\[
  2\le B<\left\lfloor\frac{\binom nk}{k+1}\right\rfloor,
  \qquad B=\lfloor2^{-128}|\F|\rfloor.
\]
The exponential alphabet is what places the explicit family in the final
regime; it is not needed for the universal endpoint equality.

For the fractional target with \(q<2^{256}\), every positive interleaving
width has the same safe radii as width one. To see this, let an interleaved
ball contain \(B+1\) distinct candidates. A linear row projection to one
symbol per site preserves membership in the scalar code and cannot increase
distance from the projected center. For each candidate pair, choose a site
where their symbols differ. Coefficient vectors that collide on that pair
lie in the proper linear hyperplane given by that site's difference.
There are \(\binom{B+1}{2}<q\) pairs, since \(B+1\le2^{128}\) and
\(B\,2^{128}\le q\). The union of their hyperplanes cannot fill the
coefficient space. One projection therefore preserves all \(B+1\)
candidates, violating the same scalar budget. A scalar violating list
embeds in one interleaving component, giving the converse.

This distinction matters in light of known capacity results.  Generic
Reed--Solomon evaluation points can, for every fixed positive gap
\(\epsilon\), achieve radius \(1-\rho-\epsilon\) over exponentially large
fields with a list size bounded in terms of \(\epsilon\)
\cite{brakensiekGopiMakam2024}. This positive-gap statement differs from the
exact endpoint equality above, and generic points do not supply the
structured smooth domain required here. Conversely, near-capacity proximity-gap
behavior can fail for particular smooth prime-field families
\cite{kambire2026}. Al Kindi \cite{kindi2026nearCapacity} constructs large
lists at vanishing slack on multiplicative subgroups, with power-of-two
lengths available at the benchmark rates.

For prescribed prime-field domains, Brakensiek et al.
\cite{brakensiekEtAl2026algorithmic} give deterministic list decoding at
fixed positive slack from capacity. Jeronimo \cite[Sections~8.2--8.4]{jeronimo2026algorithmic}
bounds fixed-width interleaved lists by a polynomial in the block length
and gives simultaneous list and MCA safety certificates at fixed slack
when the field meets the resulting budget. The exact finite threshold also
requires locating the first unsafe shell.

Recent work of Dao, Kominers, and Thaler \cite{daoKominersThaler2026}
gives quantitative list-decoding and MCA bounds on prescribed evaluation
domains, including beyond-Johnson bounds under sufficiently large
characteristic. Their positive agreement-gap counting results do not cover
the exact agreement endpoint \(A=k\); their beyond-Johnson counting
hypotheses also exclude the characteristic-three family with \(k\ge64\)
used here. Their Johnson-regime results in every characteristic remain a
separate comparison. The present binomial bounds are not asserted to improve
the sharper radius-dependent estimates in that literature, which also
includes existing Lean/ArkLib formalizations.
